\clearpage
\SourceLeaf{470}


\Needspace{9\baselineskip}\begin{center}NACHTRÄGE.\end{center}
\BookRule{0.15\linewidth}
\textit{Zu Artikel} 28. Die unbestimmte Gleichung \(ax=by\pm1\) wurde nicht zuerst vom berühmten Euler gelöst (wie dort gesagt wird), sondern schon von Bachet de Meziriac, einem Geometer des \(17\). Jahrhunderts und dem berühmten Herausgeber und Kommentator Diophants, dem der berühmte La Grange diese Ehre zugesprochen hat (\textit{Add. à l'Algèbre d'Euler p.} 525, wo zugleich die Art der Methode angegeben ist). Bachet teilte seine Entdeckung in der zweiten Auflage seines Buches \textit{Problèmes plaisans} \textit{et délectables qui se font par les nombres}, 1624, mit; in der ersten Auflage (Lyon 1612), die ich allein einsehen konnte, findet sie sich noch nicht, wird aber bereits angekündigt.

\textit{Zu den Artikeln} 151, 296, 297. Der berühmte Le Gendre hat seinen Beweis in seinem ausgezeichneten Werk \textit{Essai d'une théorie des nombres}, S. 214 ff., erneut dargelegt, jedoch ohne wesentliche Änderung; daher bleibt diese Methode auch weiterhin allen in Artikel 297 vorgebrachten Einwänden ausgesetzt. Der Satz (auf dem eine der Annahmen beruht), dass sich in jeder arithmetischen Folge \(l,\allowbreak l+k,\allowbreak l+2k\), usw., Primzahlen finden, wenn \(k\) und \(l\) keinen gemeinsamen Teiler haben, ist in diesem Werk zwar auf S. 12 ff. ausführlicher betrachtet worden; doch scheint den Anforderungen geometrischer Strenge noch nicht genügt zu sein. Aber auch dann, wenn dieser Satz vollständig bewiesen ist, bleibt die andere Annahme bestehen (dass es Primzahlen der Form \(4n+3\) gibt, für die eine gegebene, positiv genommene Primzahl der Form \(4n+1\) ein quadratischer Nichtrest ist), und ich weiß nicht, ob sie sich \textit{streng} beweisen lässt, ohne das Fundamentaltheorem selbst schon \textit{vorauszusetzen}. Ferner ist zu beachten, dass der berühmte Le Gendre diese letztere Annahme nicht stillschweigend gemacht, sondern sie auch selbst nicht verschwiegen hat, S. 221.
\SourceLeaf{471}

\textit{Zu den Artikeln} 288---293. Über denselben Gegenstand, der hier als besondere Anwendung der Theorie der ternären Formen dargestellt und hinsichtlich Strenge und Allgemeinheit so vollständig behandelt zu sein scheint, dass nichts mehr zu wünschen bleibt, hat der berühmte Le Gendre im dritten Teil seines Werkes auf S. 321---400 eine viel umfangreichere Untersuchung angestellt.\HistoricalNote{*)}{Auch ohne unseren Hinweis werden sich die Leser hüten, unsere ternären Formen mit dem zu verwechseln, was der berühmte Le Gendre \textit{forme trinaire d'un nombre} genannt hat. Mit diesem Ausdruck bezeichnete er nämlich die Zerlegung einer Zahl in drei Quadrate.} Er benutzte ganz andere Prinzipien und Methoden als wir, geriet auf diesem Wege jedoch in mehrere Schwierigkeiten, die es ihm unmöglich machten, die Hauptsätze durch einen strengen Beweis zu stützen. Diese Schwierigkeiten hat er selbst offen angegeben; doch lassen sie sich, wenn wir nicht irren, vielleicht leichter beheben als die Schwierigkeit, dass auch in dieser Untersuchung der eben erwähnte Satz (In jeder arithmetischen Folge, usw.) vorausgesetzt wird, S. 371, am Ende der Anmerkung.

\textit{Zu Artikel} 306 VIII. Im dritten Tausend negativer Determinanten wurden \(37\) irreguläre gefunden, von denen \(18\) den Irregularitätsindex \(2\) und die \(19\) übrigen den Index \(3\) haben.

\clearpage\textit{Zum selben Artikel} X. Die hier gestellte Frage konnte vor Kurzem vollständig gelöst werden. Diese Untersuchung, die mehrere Teile sowohl der höheren Arithmetik als auch der Analysis auf bemerkenswerte Weise erhellt, werden wir so bald wie möglich in der Fortsetzung dieses Werkes mitteilen. Sie hat auch gezeigt, dass der Koeffizient \(m\) in Artikel 304 \(=\gamma\pi=2{,}3458847616\) ist, wobei \(\gamma\) dieselbe Größe wie in Artikel 302 und \(\pi\), wie dort, den halben Umfang eines Kreises vom Radius \(1\) bezeichnet.
\BookRule{0.28\linewidth}


